% =====================================================================
%  guia_vancouver.tex  --  Guia de referencia de la norma Vancouver
%  Documento de trabajo del autor. NO forma parte de la tesis.
%  Compilar desde tesis/guia_vancouver/:   pdflatex guia_vancouver
%  (dos pasadas, por el indice de contenidos)
% =====================================================================
\documentclass[11pt,a4paper]{article}
\usepackage[utf8]{inputenc}\usepackage[T1]{fontenc}\usepackage{lmodern}
\usepackage[spanish,es-noquoting]{babel}
\usepackage[margin=2.3cm,top=2.6cm,bottom=2.4cm]{geometry}
\usepackage{longtable,booktabs,array,xcolor,fancyhdr,titlesec,enumitem,hyperref}
\usepackage{fancyvrb}
\usepackage[most]{tcolorbox}
\usepackage{microtype}

\definecolor{tinta}{RGB}{28,34,44}
\definecolor{gris}{RGB}{110,116,126}
\definecolor{acento}{RGB}{23,86,140}
\definecolor{fondo}{RGB}{246,247,249}
\definecolor{marco}{RGB}{206,212,220}
\definecolor{bien}{RGB}{22,110,70}
\definecolor{ojo}{RGB}{160,84,10}

\hypersetup{colorlinks=true,linkcolor=acento,urlcolor=acento,
  pdftitle={La norma Vancouver aplicada a la tesis EduFEM},
  pdfauthor={Hedy Yhassmany Oyola Sucullani}}

\pagestyle{fancy}\fancyhf{}
\fancyhead[L]{\small\color{gris} Norma Vancouver --- guía de referencia}
\fancyhead[R]{\small\color{gris}\thepage}
\renewcommand{\headrulewidth}{0.4pt}

\titleformat{\section}{\large\bfseries\color{acento}}{\thesection.}{0.6em}{}
\titlespacing*{\section}{0pt}{14pt}{6pt}
\titleformat{\subsection}{\normalsize\bfseries\color{tinta}}{\thesubsection}{0.6em}{}
\titlespacing*{\subsection}{0pt}{10pt}{4pt}
\setlist{nosep,leftmargin=1.5em}

% --- Caja de ficha por tipo de fuente ---
\newtcolorbox{ficha}[1]{
  enhanced, breakable, colback=white, colframe=marco, boxrule=0.5pt,
  arc=2pt, left=8pt, right=8pt, top=7pt, bottom=7pt,
  before skip=9pt, after skip=9pt,
  title={#1}, coltitle=white, colbacktitle=acento,
  fonttitle=\bfseries\small, attach boxed title to top left={xshift=8pt,yshift=-7pt},
  boxed title style={arc=1pt,boxrule=0pt,left=5pt,right=5pt,top=2pt,bottom=2pt}}

% --- Caja de nota al margen del texto ---
\newtcolorbox{nota}{
  enhanced, breakable, colback=fondo, colframe=fondo, boxrule=0pt,
  borderline west={2.2pt}{0pt}{acento},
  left=8pt, right=8pt, top=5pt, bottom=5pt, before skip=8pt, after skip=8pt}

% --- Rotulos internos de cada ficha ---
\newcommand{\rot}[1]{\vspace{2pt}\par{\footnotesize\bfseries\color{gris}\MakeUppercase{#1}}\par\vspace{1pt}}
\newcommand{\refv}[1]{\par{\small #1}\par}
\newcommand{\plantilla}[1]{\par{\small\itshape\color{tinta} #1}\par}
\newcommand{\anot}[1]{{\color{gris}\footnotesize #1}}

\begin{document}

% =====================================================================
\begin{center}
{\LARGE\bfseries\color{tinta} La norma Vancouver}\\[4pt]
{\large\color{tinta} Plantillas y ejemplos por tipo de fuente}\\[10pt]
{\color{gris} Tesis EduFEM --- documento de trabajo del autor}\\[10pt]
\begin{minipage}{0.9\textwidth}\small
Esta guía reúne lo que la norma Vancouver dice \emph{en sus fuentes oficiales} y lo aplica a
las 22 fuentes que la tesis cita hoy. Para cada tipo de fuente hay tres cosas: la
\textbf{plantilla} (qué elementos van y en qué orden), un \textbf{ejemplo} construido con una
fuente real de la tesis y la \textbf{entrada \texttt{.bib}} que lo produce. Al final está la
revisión de lo que la tesis imprime hoy frente a lo que pide la norma.\\[5pt]
Este documento no forma parte de la tesis.\\[5pt]
\textbf{Vancouver gobierna las citas y las referencias, no la presentación.} Desde el
2026-09-10 la tesis se presenta en formato \textbf{APA 7.ª ed.} ---márgenes, fuente,
interlineado, sangría, jerarquía de títulos, rótulos de tablas y figuras y numeración de
página---, y sigue citando en Vancouver. Los dos conviven sin chocar: APA no dice nada sobre
cómo se escribe una referencia numérica, y Vancouver no dice nada sobre el interlineado.
\end{minipage}
\end{center}
\vspace{6pt}\hrule\vspace{10pt}

{\small\tableofcontents}
\vspace{6pt}\hrule

% =====================================================================
\section{Qué es la norma Vancouver y quién la define}

En 1978, un grupo reducido de directores de revistas médicas generales se reunió de manera
informal en Vancouver (Columbia Británica, Canadá) para acordar un formato uniforme de los
manuscritos que recibían. Se los conoció como el «Grupo de Vancouver», y de ahí viene el
nombre del estilo. Su documento, los \emph{Uniform Requirements for Manuscripts Submitted to
Biomedical Journals}, se publicó ese mismo año; los formatos de las referencias
bibliográficas no los redactó el grupo, sino la \textbf{National Library of Medicine} de
Estados Unidos ---y esa división de trabajo, de 1978, es la que sigue vigente hoy---. El
grupo se institucionalizó después como \textbf{ICMJE} (International Committee of Medical
Journal Editors).

Hoy la norma vive en dos documentos, y conviene saber qué dice cada uno porque resuelven
preguntas distintas:

\begin{itemize}\setlength{\itemsep}{4pt}
\item \textbf{ICMJE}, \emph{Recommendations for the Conduct, Reporting, Editing, and
  Publication of Scholarly Work in Medical Journals} (última actualización: enero de 2026).
  Fija los \emph{principios}: cómo se numeran las citas, en qué orden, qué se puede citar y
  qué no. No da el formato de cada tipo de fuente: remite explícitamente al NLM.
\item \textbf{NLM} (National Library of Medicine, EE.\,UU.), \emph{Citing Medicine: The NLM
  Style Guide for Authors, Editors, and Publishers}, 2.ª ed., de Karen Patrias, con Dan
  Wendling como editor técnico. Es el manual completo: 26 capítulos, uno por tipo de fuente,
  con las reglas de puntuación elemento por elemento. Su resumen operativo es la página
  \emph{Sample References}, que es la que se consulta en el día a día.
\end{itemize}

\begin{nota}
\textbf{Vancouver, estilo ICMJE y estilo NLM son el mismo estilo.} Los tres nombres aparecen
en la bibliografía sobre citación y designan lo mismo; \emph{Citing Medicine} es la
encarnación actual del sistema Vancouver. Si un reglamento pide «normas Vancouver», lo que
hay que seguir es esto.
\end{nota}

\subsection*{Una advertencia que aplica a esta tesis}

Vancouver nació en biomedicina y su manual está escrito desde ahí: los ejemplos son de
revistas médicas y algunas reglas presuponen que la revista está indexada en el catálogo del
NLM. Una tesis de ingeniería civil usa el mismo esqueleto, pero hay dos puntos donde no puede
seguir la letra:

\begin{itemize}
\item Las revistas de ingeniería en general \textbf{no están en el catálogo del NLM}, de modo
  que la abreviatura del título no se puede «buscar»: hay que construirla (§\ref{sec:reglas},
  regla 5).
\item El documento está en \textbf{español}, y la norma está pensada para textos en inglés.
  Donde la norma dice \texttt{In:}, \texttt{editors}, \texttt{p.} o \texttt{2nd ed.}, un
  documento en español escribe \texttt{En:}, \texttt{editores}, \texttt{p.} y \texttt{2.ª
  ed.} Es la adaptación habitual y no es un error, siempre que se aplique de forma uniforme.
\end{itemize}

% =====================================================================
\section{Las siete reglas que gobiernan todas las referencias}
\label{sec:reglas}

Antes de las plantillas conviene fijar las reglas transversales. Casi todos los errores de
una bibliografía Vancouver son incumplimientos de alguna de estas siete.

\begin{enumerate}\setlength{\itemsep}{6pt}

\item \textbf{Numeración por orden de aparición.} Las referencias se numeran
  «consecutivamente en el orden en que se mencionan por primera vez en el texto». El número
  se asigna la primera vez y \emph{se reutiliza} cada vez que la fuente vuelve a citarse. La
  lista final va en ese orden, no alfabético.
  \anot{En tu tesis lo garantiza \texttt{sorting=none} en el preámbulo.}

\item \textbf{Autores: apellido, espacio, iniciales pegadas.} Sin coma entre apellido e
  iniciales, sin punto tras cada inicial, sin espacio entre iniciales, máximo dos iniciales
  por autor. Los autores se separan entre sí con \textbf{coma}.
  \par\vspace{3pt}\hspace{1.5em}\begin{minipage}{0.9\textwidth}\small
  \textcolor{bien}{\textbf{Bien:}} Bathe KJ \quad\textbullet\quad Lee JY, Ryu HR
  \quad\textbullet\quad Timoshenko SP, Goodier JN\\
  \textcolor{ojo}{\textbf{Mal:}} Bathe, K.-J. \quad\textbullet\quad Lee J. Y. y Ryu H.-R.
  \end{minipage}

\item \textbf{Cuántos autores se listan.} Acá hay una diferencia real entre los dos documentos
  oficiales, y hay que elegir uno:
  \par\vspace{3pt}
  \begin{minipage}{0.95\textwidth}\small
  \begin{itemize}
  \item \emph{Citing Medicine} dice: «Give all authors, regardless of the number» ---todos,
    sin importar cuántos.
  \item La página \emph{Sample References} del NLM (la convención ICMJE, la más difundida)
    dice: «List the first six authors followed by \emph{et al.}»
  \end{itemize}
  \end{minipage}
  \par\vspace{3pt}
  \anot{Tu tesis usa la segunda: seis y \emph{et al.} (\texttt{maxbibnames=6},
  \texttt{minbibnames=6}). Es la correcta para NumPy, que tiene 26 autores, y SciPy, que
  tiene 35.}

\item \textbf{Mayúsculas del título: sólo la primera palabra.} La norma es explícita:
  «capitalize only the first word of a title, proper nouns, proper adjectives, acronyms, and
  initialisms». Es decir, el título va como una oración, no en \emph{title case} inglés,
  aunque en la portada del libro aparezca con todas las iniciales en mayúscula.
  \par\vspace{3pt}\hspace{1.5em}\begin{minipage}{0.9\textwidth}\small
  \textcolor{bien}{\textbf{Bien:}} The finite element method: its basis and fundamentals\\
  \textcolor{ojo}{\textbf{Mal:}} The Finite Element Method: Its Basis and Fundamentals
  \end{minipage}
  \par\vspace{3pt}
  \anot{Las siglas y nombres propios se conservan: NumPy, SciPy, Python, ED-Elas2D, SAP2000.}

\item \textbf{El título de la revista va abreviado.} Se abrevian las palabras significativas
  y se omiten artículos, conjunciones y preposiciones. La abreviatura oficial se busca en el
  \emph{NLM Catalog}; si la revista no está ahí ---el caso de casi toda la ingeniería--- se
  construye con la norma ISO 4, usando la \emph{List of Title Word Abbreviations} (LTWA) del
  Centro Internacional del ISSN. En estilo NLM la abreviatura va \textbf{sin puntos}.
  \par\vspace{3pt}\hspace{1.5em}\begin{minipage}{0.9\textwidth}\small
  Computer Applications in Engineering Education $\rightarrow$ \textbf{Comput Appl Eng Educ}\\
  European Journal of Engineering Education $\rightarrow$ \textbf{Eur J Eng Educ}\\
  Nature Methods $\rightarrow$ \textbf{Nat Methods} \quad\textbullet\quad
  Nature $\rightarrow$ \textbf{Nature} \anot{(de una sola palabra: no se abrevia)}
  \end{minipage}

\item \textbf{Rangos de página comprimidos.} No se repiten los dígitos que ya aparecieron en
  el número inicial.
  \par\vspace{3pt}\hspace{1.5em}\begin{minipage}{0.9\textwidth}\small
  1159\,-\,1173 $\rightarrow$ \textbf{1159-73} \quad\textbullet\quad
  872\,-\,886 $\rightarrow$ \textbf{872-86} \quad\textbullet\quad
  157\,-\,169 $\rightarrow$ \textbf{157-69} \quad\textbullet\quad
  1007\,-\,1027 $\rightarrow$ \textbf{1007-27}
  \end{minipage}

\item \textbf{La puntuación del pie de imprenta es fija:} \texttt{Lugar: Editorial; año.}
  Dos puntos después del lugar, \textbf{punto y coma} antes del año, punto final. Las
  ciudades de EE.\,UU. y Canadá llevan la abreviatura de dos letras del estado o provincia
  entre paréntesis; las demás llevan el país sólo si hace falta para desambiguar.
  \par\vspace{3pt}\hspace{1.5em}\begin{minipage}{0.9\textwidth}\small
  Upper Saddle River (NJ): Prentice Hall; 1996. \quad\textbullet\quad
  Oxford: Elsevier Butterworth-Heinemann; 2005.
  \end{minipage}
\end{enumerate}

\subsection*{Dos reglas más que vas a necesitar}

\textbf{Títulos que no están en inglés.} Si el título está en alfabeto latino ---el caso del
español--- se \textbf{deja tal cual}, en su idioma original. La traducción entre corchetes es
para lenguas que hay que romanizar o traducir (chino, japonés, ruso, árabe) o para cuando el
documento que escribís está en inglés. En una tesis en español, tus fuentes en español van sin
corchetes y tus fuentes en inglés también van sin traducir.

\textbf{Datos que no existen.} Cuando un ejemplar no consigna lugar, editorial o año, la norma
no permite omitir el elemento en silencio: se declara la ausencia entre corchetes
---\texttt{[place unknown]}, \texttt{[publisher unknown]}, \texttt{[date unknown]}---, que en
un documento en español se escribe \texttt{[lugar desconocido]}, \texttt{[editorial
desconocida]}, \texttt{[fecha desconocida]}. Esto te aplica directamente en la monografía de
Álvarez de Zayas.

% =====================================================================
\section{Los cuatro tipos que usás hoy}

Las 22 fuentes de \texttt{referencias.bib} se reparten en cinco tipos: 11 libros, 7 artículos
de revista, 2 informes técnicos, 1 manual de organización y 1 monografía inédita. Una ficha
por tipo.

% ---------------------------------------------------------------------
\begin{ficha}{A \textperiodcentered\ LIBRO COMPLETO \hfill \normalfont\texttt{@book} \textperiodcentered\ 11 entradas}
\rot{plantilla}
\plantilla{Autor AA, Autor BB, Autor CC. Título del libro: subtítulo. Edición. Lugar
(Estado): Editorial; año.}
\anot{La edición se omite si es la primera. La paginación total (\texttt{245 p.}) es opcional
y en una tesis no se pone.}

\rot{ejemplo --- tu fuente [1]}
\refv{Zienkiewicz OC, Taylor RL, Zhu JZ. The finite element method: its basis and
fundamentals. 6.ª ed. Oxford: Elsevier Butterworth-Heinemann; 2005.}

\rot{ejemplo --- tu fuente en español}
\refv{García-Córdoba F. La investigación tecnológica: investigar, idear e innovar en
ingenierías y ciencias sociales. México (DF): Limusa; 2005.}

\rot{entrada .bib}
\begin{Verbatim}[fontsize=\small,xleftmargin=4pt]
@book{zienkiewicz2013fem,
  author    = {Zienkiewicz, O. C. and Taylor, R. L. and Zhu, J. Z.},
  title     = {The Finite Element Method: Its Basis and Fundamentals},
  edition   = {6},
  publisher = {Elsevier Butterworth-Heinemann},
  location  = {Oxford},
  year      = {2005},
}
\end{Verbatim}
\anot{Campos mínimos: \texttt{author}, \texttt{title}, \texttt{publisher}, \texttt{location},
\texttt{year}. \texttt{location}, no \texttt{address}. El \texttt{isbn} es opcional y sólo si
lo leíste de la página de créditos.}
\end{ficha}

% ---------------------------------------------------------------------
\begin{ficha}{B \textperiodcentered\ ARTÍCULO DE REVISTA \hfill \normalfont\texttt{@article} \textperiodcentered\ 7 entradas}
\rot{plantilla}
\plantilla{Autor AA, Autor BB. Título del artículo. Abrev Revista.
año;volumen(fascículo):pág-pág. doi:10.xxxx/yyyy}
\anot{Sin comillas en el título del artículo, sin la palabra «En:» y sin coma antes del año.
Todo el bloque final va pegado: punto y coma, paréntesis, dos puntos.}

\rot{ejemplo --- dos autores}
\refv{Lee JY, Ryu HR. Interactive simulation of eigenmodes for instruction of finite element
behavior. Comput Appl Eng Educ. 2015;23(6):872-86. doi:10.1002/cae.21659}

\rot{ejemplo --- más de seis autores}
\refv{Harris CR, Millman KJ, van der Walt SJ, Gommers R, Virtanen P, Cournapeau D, et al.
Array programming with NumPy. Nature. 2020;585(7825):357-62.
doi:10.1038/s41586-020-2649-2}

\rot{ejemplo --- fascículo no impreso en el ejemplar}
\refv{Pérez-Santiago R, Campos E. FEM education in undergraduate studies: industry-informed
research. Comput Appl Eng Educ. 2023;31(5):1159-73. doi:10.1002/cae.22627}
\anot{El ejemplar no imprime el fascículo; el «5» sale de la marca de agua de Wiley y de la
ficha de CrossRef del DOI (verificada el 2026-09-11). Si no puede comprobarse, se omite:
no se inventa.}

\rot{entrada .bib}
\begin{Verbatim}[fontsize=\small,xleftmargin=4pt]
@article{lee2015eigenmodes,
  author       = {Lee, Jae Young and Ryu, Hee-Ryong},
  title        = {Interactive Simulation of Eigenmodes for Instruction
                  of Finite Element Behavior},
  journaltitle = {Computer Applications in Engineering Education},
  volume       = {23},
  number       = {6},
  pages        = {872--886},
  year         = {2015},
  doi          = {10.1002/cae.21659},
}
\end{Verbatim}
\anot{\texttt{journaltitle}, no \texttt{journal}. Las páginas se escriben completas en el
\texttt{.bib} (con doble guion); comprimirlas es tarea del estilo, no tuya.}
\end{ficha}

% ---------------------------------------------------------------------
\begin{ficha}{C \textperiodcentered\ INFORME CIENTÍFICO O TÉCNICO \hfill \normalfont\texttt{@techreport} \textperiodcentered\ 2 entradas}
\rot{plantilla}
\plantilla{Autor AA, Autor BB. Título del informe. Lugar (Estado): Institución; año. Informe
n.º: XXXX.}
\anot{En inglés el rótulo es \texttt{Report No.:}; si hay contrato o subvención se agregan
después como \texttt{Contract No.:} o \texttt{Grant No.:}. La institución ocupa el lugar de la
editorial.}

\rot{ejemplo --- tu fuente de MMS}
\refv{Salari K, Knupp P. Code verification by the method of manufactured solutions.
Albuquerque (NM): Sandia National Laboratories; 2000. Informe n.º: SAND2000-1444.}

\rot{ejemplo --- cinco autores}
\refv{Stimpson CJ, Ernst CD, Knupp P, Pébay PP, Thompson D. The Verdict geometric quality
library. Albuquerque (NM): Sandia National Laboratories; 2007. Informe n.º: SAND2007-1751.}

\rot{entrada .bib}
\begin{Verbatim}[fontsize=\small,xleftmargin=4pt]
@techreport{salari2000mms,
  author      = {Salari, Kambiz and Knupp, Patrick},
  title       = {Code Verification by the Method of Manufactured Solutions},
  institution = {Sandia National Laboratories},
  location    = {Albuquerque, NM},
  number      = {SAND2000-1444},
  year        = {2000},
}
\end{Verbatim}
\anot{El número del informe va en \texttt{number}: es el dato que permite recuperarlo del
repositorio de Sandia, y sin él la referencia no identifica el documento.}
\end{ficha}

% ---------------------------------------------------------------------
\begin{ficha}{D \textperiodcentered\ MANUAL U OBRA DE UNA ORGANIZACIÓN \hfill \normalfont\texttt{@manual} \textperiodcentered\ 1 entrada}
\rot{plantilla}
\plantilla{Organización. Título del manual. Versión o revisión. Lugar (Estado): Organización;
año.}
\anot{Cuando no hay autor personal, \textbf{la organización ocupa el lugar del autor} y se
escribe completa, sin invertir nada. Si además es la editora, la norma permite abreviarla en
la posición de editorial: «Berkeley (CA): The Company;».}

\rot{ejemplo --- tu manual de SAP2000}
\refv{Computers \& Structures, Inc. CSI analysis reference manual for SAP2000, ETABS, SAFE and
CSiBridge. Rev. 18. [Estados Unidos]: Computers \& Structures, Inc.; 2017.}
\anot{El ejemplar no consigna ciudad en ningún lado: la portada dice sólo «Proudly developed
in the United States of America» y la página de copyright da la empresa y su web, sin
domicilio. El lugar va \textbf{inferido de la propia publicación y entre corchetes}, que es
la convención del NLM para un dato que no está impreso.}

\rot{entrada .bib}
\begin{Verbatim}[fontsize=\small,xleftmargin=4pt]
@manual{csi2017sap2000,
  author    = {{Computers \& Structures, Inc.}},
  title     = {{CSI} Analysis Reference Manual for {SAP2000}, {ETABS},
               {SAFE} and {CSiBridge}},
  publisher = {Computers \& Structures, Inc.},
  location  = {[Estados Unidos]},
  year      = {2017},
  note      = {Rev. 18},
}
\end{Verbatim}
\anot{Las llaves dobles \texttt{\{\{...\}\}} son obligatorias: sin ellas biber lee «Inc.» como
apellido y «Computers \& Structures» como nombre de pila. La organización va dos veces, en
\texttt{author} y en \texttt{publisher}: es lo que pide el NLM cuando una entidad publica su
propia obra. No uses \texttt{organization} para eso: biblatex 3.21 lo imprime \emph{antes} del
pie de imprenta y la referencia sale invertida («Computers \& Structures, Inc. [Estados
Unidos], 2017»); sólo \texttt{publisher} ocupa el lugar «Lugar: Editorial, año». La revisión
va en \texttt{note} y no en \texttt{version}, que imprimiría «Ver. Rev. 18».}
\end{ficha}

% ---------------------------------------------------------------------
\begin{ficha}{E \textperiodcentered\ MONOGRAFÍA INÉDITA \hfill \normalfont\texttt{@unpublished} \textperiodcentered\ 1 entrada}
\rot{plantilla}
\plantilla{Autor AA. Título. Nota de procedencia. [lugar desconocido]: [editorial
desconocida]; [fecha desconocida].}
\anot{Vancouver no permite omitir en silencio lugar, editorial o fecha: se declara la
ausencia entre corchetes. La nota de procedencia es lo que vuelve localizable un documento
sin pie de imprenta: qué es, dónde circula y qué fecha lleva el archivo.}

\rot{ejemplo --- tu monografía de Álvarez de Zayas}
\refv{Álvarez de Zayas C. Metodología de la investigación científica. Monografía inédita de
80 páginas, en formato electrónico (PDF fechado en 2007 según sus metadatos), empleada en la
Carrera de Ingeniería Civil de la UATF como referencia del marco metodológico. [lugar
desconocido]: [editorial desconocida]; [fecha desconocida].}

\rot{entrada .bib}
\begin{Verbatim}[fontsize=\small,xleftmargin=4pt]
@unpublished{alvarez_metodologia,
  author = {{\'A}lvarez de Zayas, Carlos},
  title  = {Metodolog{\'i}a de la investigaci{\'o}n cient{\'i}fica},
  note   = {Monograf{\'i}a in{\'e}dita de 80 p{\'a}ginas, ... [lugar desconocido]:
            [editorial desconocida]; [fecha desconocida]},
}
\end{Verbatim}
\anot{La fecha desconocida no puede ir en \texttt{year}: biber la parsea como fecha y la
descarta con aviso. Va al final de \texttt{note}, que es lo que la deja en el renglón
correcto. El estilo imprime el título entre comillas, como hace con todo
\texttt{@unpublished}: es una limitación de \texttt{numeric-comp}, la misma de los
artículos.}
\end{ficha}

% =====================================================================
\section{Los tipos que todavía no usás pero podrías necesitar}

Si en la defensa o en una revisión aparece una fuente de otro tipo, estas son las cinco
plantillas que cubren lo que razonablemente puede salir en esta tesis. Los ejemplos marcados
\anot{[NLM]} son los ejemplos oficiales del manual, reproducidos tal cual para que sirvan de
patrón de puntuación.

\subsection*{Capítulo o parte de un libro \hfill \normalfont\small\texttt{@incollection}}
\plantilla{Autor del capítulo. Título del capítulo. En: Editor AA, Editor BB, editores.
Título del libro. Edición. Lugar: Editorial; año. p. inicio-fin.}
\refv{\anot{[NLM]} Meltzer PS, Kallioniemi A, Trent JM. Chromosome alterations in human solid
tumors. In: Vogelstein B, Kinzler KW, editors. The genetic basis of human cancer. New York:
McGraw-Hill; 2002. p. 93-113.}
\anot{El único caso en que Vancouver escribe \texttt{p.} antes de las páginas es este y el de
las actas de congreso; en los artículos de revista no lleva rótulo.}

\subsection*{Ponencia publicada en actas de congreso \hfill \normalfont\small\texttt{@inproceedings}}
\plantilla{Autor AA, Autor BB. Título de la ponencia. En: Editor AA, editor. Título de las
actas; año mes días; ciudad, país. Lugar: Editorial; año. p. inicio-fin.}
\refv{\anot{[NLM]} Christensen S, Oppacher F. An analysis of Koza's computational effort
statistic for genetic programming. In: Foster JA, Lutton E, Miller J, Ryan C, Tettamanzi AG,
editors. Genetic programming. EuroGP 2002: Proceedings of the 5th European Conference on
Genetic Programming; 2002 Apr 3-5; Kinsdale, Ireland. Berlin: Springer; 2002. p. 182-91.}
\anot{Fijate que la fecha del congreso y la de publicación de las actas son dos datos
distintos y van los dos.}

\subsection*{Tesis \hfill \normalfont\small\texttt{@thesis}}
\plantilla{Autor AA. Título de la tesis [tesis de licenciatura]. Lugar (Estado):
Universidad; año.}
\refv{\anot{[NLM]} Borkowski MM. Infant sleep and feeding: a telephone survey of Hispanic
Americans [dissertation]. Mount Pleasant (MI): Central Michigan University; 2002.}
\anot{El tipo de trabajo va entre corchetes después del título: \texttt{[tesis de
licenciatura]}, \texttt{[tesis de maestría]}, \texttt{[tesis doctoral]}.}

\subsection*{Documento o sitio web \hfill \normalfont\small\texttt{@online}}
\plantilla{Autor u organización. Título [Internet]. Lugar: Editor; año [citado año mes día].
Disponible en: URL}
\refv{Oyola Sucullani HY. EduFEM: software educativo de elementos finitos [Internet]. GitHub;
2026 [citado 2026 sep 10]. Disponible en: https://github.com/Sucullani/SoftwareED}
\anot{Los dos elementos que distinguen una fuente en línea son el designador
\texttt{[Internet]} tras el título y la \textbf{fecha de consulta} entre corchetes: sin ella la
referencia no dice a qué versión del sitio te referís. Hoy tu tesis pone el repositorio en
nota al pie, que es una alternativa legítima: la nota no entra en la numeración de
referencias.}

\subsection*{Programa de computadora \hfill \normalfont\small\texttt{@software}}
\plantilla{Autor u organización. Nombre: descripción. Versión X.Y [software]. Fecha [citado
año mes día]. Disponible en: URL}
\refv{\anot{[NLM]} Hayes B, Tesar B, Zurow K. OTSoft: Optimality Theory Software. Version
2.3.2 [software]. 2013 Jan 14 [cited 2015 Feb 14]. Available from:
https://linguistics.ucla.edu/people/hayes/otsoft/}

\begin{nota}
\textbf{NumPy y SciPy están bien citados como artículos, no como software.} Ambos proyectos
publicaron un artículo en \emph{Nature} y \emph{Nature Methods} respectivamente y piden
explícitamente que se cite ese artículo. Citar la biblioteca como programa sería
técnicamente admisible pero peor: el artículo es la fuente citable, con DOI permanente y
autoría acreditada.
\end{nota}

% =====================================================================
\section{Tus 22 fuentes en Vancouver}

Inventario de lo que hay en \texttt{referencias.bib}, agrupado por tipo. Sirve como control:
cualquier fuente nueva tiene que caer en una de estas cuatro filas o abrir una fila nueva con
su plantilla de la §4.

\vspace{4pt}
{\small
\begin{longtable}{@{}p{3.3cm}p{1.0cm}p{11.0cm}@{}}
\toprule
\textbf{Tipo} & \textbf{N.º} & \textbf{Fuentes} \\
\midrule\endhead
Libro \newline \anot{\texttt{@book}} & 11 &
Zienkiewicz \textperiodcentered\ Bathe \textperiodcentered\ Cook \textperiodcentered\ Hughes
\textperiodcentered\ Reddy \textperiodcentered\ Oñate \textperiodcentered\ Timoshenko
\textperiodcentered\ Strang \textperiodcentered\ Oberkampf
\textperiodcentered\ Hernández-Sampieri \textperiodcentered\ García-Córdoba \\
\addlinespace[3pt]
Monografía inédita \newline \anot{\texttt{@unpublished}} & 1 & Álvarez de Zayas \\
\addlinespace[3pt]
Artículo \newline \anot{\texttt{@article}} & 7 &
Harris (NumPy) \textperiodcentered\ Virtanen (SciPy) \textperiodcentered\ Bishay
\textperiodcentered\ Lee 2015a \textperiodcentered\ Lee y Ryu 2015b \textperiodcentered\
Suárez (ED-Elas2D) \textperiodcentered\ Pérez-Santiago y Campos \\
\addlinespace[3pt]
Informe técnico \newline \anot{\texttt{@techreport}} & 2 &
Salari y Knupp (MMS, SAND2000-1444) \textperiodcentered\ Stimpson (Verdict, SAND2007-1751) \\
\addlinespace[3pt]
Manual \newline \anot{\texttt{@manual}} & 1 &
CSI, manual de referencia de SAP2000 \\
\bottomrule
\end{longtable}}

\begin{nota}
\textbf{Las tres revistas que citás y su abreviatura NLM/ISO 4:} \emph{Computer Applications
in Engineering Education} $\rightarrow$ Comput Appl Eng Educ (cuatro artículos);
\emph{European Journal of Engineering Education} $\rightarrow$ Eur J Eng Educ (uno);
\emph{Nature} y \emph{Nature Methods} $\rightarrow$ Nature y Nat Methods (uno cada una). Las
dos últimas son las únicas de tu bibliografía que figuran en el catálogo del NLM; las dos
primeras hay que abreviarlas con la LTWA.
\end{nota}

% =====================================================================
\section{Lo que imprime tu tesis hoy frente a lo que pide la norma}

Tu preámbulo no carga un estilo Vancouver: arma uno aproximado sobre el estilo
\texttt{numeric-comp} de \texttt{biblatex}, con \texttt{sorting=none} y ajustes en el formato
de los nombres.
Eso resuelve lo esencial ---numeración por orden de aparición y apellido con iniciales--- pero
deja diferencias en la puntuación. La tabla separa lo que es \textbf{error corregible} de lo
que es \textbf{decisión ya tomada} y documentada en \texttt{tesis/README.md}.

\vspace{4pt}
{\small
\begin{longtable}{@{}p{2.5cm}p{4.3cm}p{4.3cm}p{3.9cm}@{}}
\toprule
\textbf{Elemento} & \textbf{Lo que imprime hoy} & \textbf{Lo que pide Vancouver} &
\textbf{Veredicto} \\
\midrule\endhead

Iniciales con guion &
\texttt{Bathe KJ} \newline \texttt{Ryu HR} &
\texttt{Bathe KJ} \newline \texttt{Ryu HR} &
\textcolor{bien}{\textbf{Corregido}} (§7). Salía \texttt{K.-J.} \\
\addlinespace[3pt]

Lugar de edición &
Las 15 entradas que lo exigen lo traen &
El lugar es obligatorio en libros, informes y manuales &
\textcolor{bien}{\textbf{Corregido}} (§7). Faltaba en tres. \\
\addlinespace[3pt]

Editorial del manual &
\texttt{...: Computers \& Structures, Inc.; 2017} &
La organización va de autor \emph{y} de editorial (\texttt{publisher}) &
\textcolor{bien}{\textbf{Corregido}} (§7, ítems 3 y 4). Salía sin editorial y luego invertida. \\
\addlinespace[3pt]

Título del artículo &
Entre comillas: \newline \guillemotleft Interactive Simulation\ldots\guillemotright &
Sin comillas &
Limitación de \texttt{numeric-comp}. \\
\addlinespace[3pt]

Enlace a la revista &
\texttt{En: Computer Applications\ldots} &
La revista sigue directo, sin \texttt{En:} &
Limitación de \texttt{numeric-comp}. \\
\addlinespace[3pt]

Bloque final del artículo &
\texttt{23.2 (2015), págs. 157-169} &
\texttt{2015;23(2):157-69} &
Limitación de \texttt{numeric-comp}. \\
\addlinespace[3pt]

Editorial y año &
\texttt{Oxford: Elsevier\ldots, 2005} &
\texttt{Oxford: Elsevier\ldots; 2005} &
Limitación de \texttt{numeric-comp}. \\
\addlinespace[3pt]

Nombre de la revista &
Completo, sin abreviar &
Abreviado (Comput Appl Eng Educ) &
\textcolor{bien}{\textbf{Decisión tomada}} por legibilidad en un texto en español. \\
\addlinespace[3pt]

Último separador de autores &
\texttt{Taylor RL y Zhu JZ} &
\texttt{Taylor RL, Zhu JZ} &
\textcolor{bien}{\textbf{Decisión tomada}}: es la convención de babel-spanish. \\
\addlinespace[3pt]

Mayúsculas del título &
\emph{Title Case} inglés, como en la portada &
Sólo la primera palabra y los nombres propios &
Decisión del autor. Es lo único de esta tabla que se arregla editando el \texttt{.bib}. \\
\addlinespace[3pt]

Clave y año &
\texttt{zienkiewicz2013fem} para un libro de 2005 &
--- &
\textcolor{bien}{\textbf{Sin efecto.}} Con estilo numérico el lector nunca ve la clave. \\
\bottomrule
\end{longtable}}

\begin{nota}
\textbf{Las cuatro filas de «limitación de \texttt{numeric-comp}» tienen una única causa y una
única solución posible}: existe en CTAN el paquete \texttt{biblatex-vancouver}, que implementa
el estilo completo y se activa con \texttt{style=vancouver}. Cambiarlo reescribiría de golpe
la puntuación de las 22 referencias ---y también revertiría las dos decisiones ya tomadas
(abreviaría las revistas y pondría coma en vez de «y»)---. Es una decisión de una sola vez, no
un arreglo parcial: o se adopta el estilo entero o se deja el actual documentado como
«numérico secuencial tipo Vancouver», que es lo que hoy dice el \texttt{README}.
\end{nota}

% =====================================================================
\section{Las correcciones aplicadas (2026-09-10)}

\subsection*{1. Iniciales de nombres compuestos}

\texttt{biblatex} define \texttt{\textbackslash bibinithyphendelim} como \texttt{.-}, y por eso
Klaus-Jürgen salía \texttt{K.-J.} y Hee-Ryong salía \texttt{H.-R.} en vez de \texttt{KJ} y
\texttt{HR}. Los tres \texttt{\textbackslash renewcommand} del preámbulo
(\texttt{bibinitperiod}, \texttt{bibinitdelim} y el que faltaba) los cubre de una sola vez la
opción \texttt{terseinits}, que es lo que ahora está puesto:

\begin{Verbatim}[fontsize=\small,xleftmargin=4pt]
\usepackage[backend=biber,
            style=numeric-comp,
            sorting=none,
            terseinits=true,     % <-- reemplaza los tres \renewcommand
            ...
\end{Verbatim}
\anot{La opción fija \texttt{bibinitperiod}, \texttt{bibinitdelim} y
\texttt{bibinithyphendelim} en vacío, que es exactamente el formato Vancouver.}

\subsection*{2. Lugares de edición faltantes}

Faltaba \texttt{location} en tres entradas, y la norma lo exige. Cada dato salió del ejemplar,
no de la web:

\begin{itemize}\setlength{\itemsep}{3pt}
\item \texttt{cook2002concepts} $\rightarrow$ \texttt{Hoboken (NJ)}. La portada sólo dice
  «John Wiley \& Sons, Inc.»; la ciudad está en la página de créditos, en el domicilio del
  editor («111 River Street, Hoboken, NJ 07030»).
\item \texttt{csi2017sap2000} $\rightarrow$ \texttt{[Estados Unidos]}. No hay ciudad en
  ninguna parte del manual: se infiere el país de la portada y va entre corchetes.
\item \texttt{alvarez\_metodologia} $\rightarrow$ \texttt{[lugar desconocido]}, más
  \texttt{[editorial desconocida]} y \texttt{[fecha desconocida]}. El ejemplar no trae pie de
  imprenta, y la norma no deja omitir en silencio un elemento obligatorio.
\end{itemize}

\anot{La fecha desconocida no puede ir en \texttt{year}: biber la parsea como fecha y la
descarta con aviso. Va pegada a \texttt{publisher}, que es lo que la deja en el renglón
correcto.}

\subsection*{3. El manual salía sin editorial}

En \texttt{@manual}, biblatex no imprime \texttt{author} en el lugar de la editorial. Con sólo
\texttt{author} la referencia terminaba en «\texttt{[Estados Unidos], 2017}», sin editorial. La
organización va ahora \textbf{dos veces}, en \texttt{author} y en \texttt{publisher}: es lo que
pide el NLM cuando una entidad publica su propia obra («Chicago: The Association;»). (El
2026-09-10 se había puesto en \texttt{organization}, y biblatex 3.21 imprime ese campo antes
del pie de imprenta: la referencia salió invertida; corregido el 2026-09-16, ítem 4.)
También se sacó \texttt{version = \{Rev. 18\}}, porque biblatex le antepone «Ver.» e imprimía
«\texttt{Ver. Rev. 18}»; la revisión vive ahora en \texttt{note}.

\subsection*{4. Correcciones del 2026-09-16 (auditoría por sesiones, H-42, H-44 y H-45)}

\begin{itemize}\setlength{\itemsep}{3pt}
\item \textbf{El manual salía invertido.} Con la organización en \texttt{organization}, la
  referencia imprimía «Computers \& Structures, Inc. [Estados Unidos], 2017»: biblatex 3.21
  coloca ese campo \emph{antes} del pie de imprenta. Pasó a \texttt{publisher}, que es el
  que ocupa el lugar «Lugar: Editorial, año». La ficha D y la corrección 3 quedaron
  actualizadas.
\item \textbf{Reddy sin nota de edición.} Se quitó \texttt{note = \{International Edition\}},
  por la misma regla con que se había quitado la nota de Bathe: Vancouver no contempla notas
  entre el título y el pie de imprenta.
\item \textbf{Oñate con un solo editor.} «Barcelona: CIMNE» en vez de «CIMNE y Springer»: el
  NLM manda consignar un único editor —el primero o el más destacado— y no listar varios.
\item \textbf{García-Córdoba} pasa a «México (DF)», como el ejemplo de la ficha A.
\item \textbf{Bishay} deja anotado en el \texttt{.bib} de dónde salen volumen, fascículo y
  páginas (ficha de CrossRef del DOI, verificada el 2026-09-11): el ejemplar es la versión
  \emph{Early View}, sin esos datos impresos.
\item \textbf{Álvarez de Zayas} pasa a la ficha E, con nota de procedencia; la guía ya no lo
  cuenta entre los libros.
\item \textbf{Localizadores de página.} La regla es ahora una sola y está declarada en la
  Introducción de la tesis: llevan localizador todas las citas que sostienen una ecuación, un
  valor numérico, un umbral, una definición o una atribución concreta, y no lo llevan las
  citas de encuadre. Las citas múltiples con localizador (\texttt{\textbackslash autocites}) se separan con punto
  y coma —delimitador de \texttt{\textbackslash parencites}, redeclarado en el preámbulo— para que
  cada página quede junto a su fuente; dentro de una cita simple de varias claves sigue la coma.
\end{itemize}

% =====================================================================
\section{Checklist antes de entregar}

\begin{enumerate}\setlength{\itemsep}{3pt}
\item ¿Cada entrada del \texttt{.bib} tiene los campos mínimos de su ficha (§3)?
\item ¿Todos los libros y manuales tienen \texttt{location}?
\item ¿Hay algún DOI, ISBN o número de página que no hayas leído del ejemplar? Si no lo
  confirmaste, el campo se borra: una entrada corta y correcta vale más que una completa e
  inventada.
\item ¿Corriste \texttt{biber} después de tocar el \texttt{.bib}? Si una cita sale como
  \texttt{[?]}, es eso.
\item ¿Toda entrada del \texttt{.bib} se cita al menos una vez? \texttt{biblatex} no imprime
  las no citadas, así que una entrada huérfana no da error: desaparece en silencio.
\item ¿Las citas nuevas quedaron registradas en \texttt{respaldo\_citas.tex} con su página?
\end{enumerate}

% =====================================================================
\section{Fuentes de esta guía}

Todo lo que dice este documento sale de los textos oficiales de la norma, no de guías de
terceros:

\begin{enumerate}\setlength{\itemsep}{3pt}
\item International Committee of Medical Journal Editors. Recommendations for the conduct,
  reporting, editing, and publication of scholarly work in medical journals [Internet]. ICMJE;
  act. enero de 2026 [citado 2026 sep 10]. Disponible en:
  \url{https://www.icmje.org/recommendations/}
\item Patrias K. Citing medicine: the NLM style guide for authors, editors, and publishers
  [Internet]. 2.ª ed. Wendling DL, editor técnico. Bethesda (MD): National Library of
  Medicine (US); 2007- [act. 2 oct. 2015; citado 2026 sep 10]. Disponible en:
  \url{https://www.ncbi.nlm.nih.gov/books/NBK7256/}
\item National Library of Medicine. Samples of formatted references for authors of journal
  articles [Internet]. Bethesda (MD): NLM [citado 2026 sep 10]. Disponible en:
  \url{https://www.nlm.nih.gov/bsd/uniform_requirements.html}
\item NLM Catalog: journals referenced in the NCBI databases [Internet]. Bethesda (MD): NLM
  [citado 2026 sep 10]. Disponible en: \url{https://www.ncbi.nlm.nih.gov/nlmcatalog/journals}
  \anot{--- para buscar la abreviatura oficial de una revista.}
\item ISSN International Centre. List of title word abbreviations (LTWA) [Internet] [citado
  2026 sep 10]. Disponible en: \url{https://www.issn.org/services/online-services/access-to-the-ltwa/}
  \anot{--- para construir la abreviatura de una revista que no está en el catálogo del NLM.}
\end{enumerate}

\vspace{8pt}\hrule\vspace{4pt}
\anot{Las cinco referencias de arriba están escritas en Vancouver, con las plantillas de este
mismo documento, a modo de ejemplo final.}

\end{document}
